\section{Method}
\label{sec:method}

We use the continuous bitstream diffusion model of \citet{method:cobit} as the generative core.
Its Gaussian corruption, probability-space objective, entropy-derived schedule, self-conditioning,
sampler and transformer backbone are inherited. Our contribution is the interface that makes this
core multimodal: a shared binary representation, a fixed sequence layout and a partial-denoising
objective that selects the task without changing the architecture. Figure and implementation
details of the inherited components are deferred to Appendix~\ref{app:method}.

\subsection{CoBit generative core}
\label{sec:method:cobit}

CoBit writes each token from a vocabulary of size $V$ as an
$m=\lceil\log_2 V\rceil$-bit code. A sequence of $N$ tokens becomes
$\vx_0\in\{0,1\}^{S}$, where $S=Nm$, and is corrupted as
$\vx_\sigma=\vx_0+\sigma\bm{\epsilon}$ with
$\bm{\epsilon}\sim\mathcal N(\vzero,\mathbf I)$. The denoiser
$D_\theta(\vx_\sigma,\sigma)\in(0,1)^S$ predicts the probability that each clean bit is one and is
trained with the EDM-weighted probability-space loss \citep{image:edm}
\begin{equation}
\Ls(\theta)=\E\!\left[w(\sigma)\,\frac{1}{S}
\left\|D_\theta(\vx_\sigma,\sigma)-\vx_0\right\|^2\right],
\qquad
w(\sigma)=\frac{\sigma^2+\sigma_{\mathrm{data}}^2}{\sigma^2\sigma_{\mathrm{data}}^2},
\label{eq:loss}
\end{equation}
with $\sigma_{\mathrm{data}}=1/2$. An analytic single-bit posterior is added to a learned contextual residual, self-conditioning feeds
the previous clean estimate back to the model, and an online entropy-rate estimate allocates both
training noise levels and sampling steps. Sampling uses a first-order probability-flow solver. We
use one input projection, one transformer trunk and one bitwise output head; the $m$ bits at a
position are patched into one trunk token so attention runs over $N$ positions rather than $Nm$
bits. These components follow CoBit; Appendix~\ref{app:method} gives the settings needed for
reproduction.

\subsection{One representation for several modalities}
\label{sec:method:multimodal}

Frozen tokenizers map each modality to integer codes, which are placed in disjoint ranges of one
shared code space together with padding and structural markers. The common bit width covers the
union of these ranges, so every position has the same representation and passes through the same
projection, backbone and bitwise head. Tokenizers and decoders remain modality-specific components
outside the denoiser; inside it there are no modality-specific embeddings, heads, losses or
training stages.

Each example follows a fixed sequence template with one segment per modality. Segment markers stay
clean, while content positions are truncated or padded to fixed budgets. Two-dimensional grids are
flattened in a fixed order\ifdefined\toyimageonly{} (row by row for the raw digit images of
\S\ref{sec:exp:imagetext})\else: row by row for raw digit images, and along a locality-preserving
(Hilbert) curve for quantised natural-image codes\fi. Aligned modalities may share a position. In the protein model, each residue
is one 18-bit position containing a 5-bit amino-acid code and the 13 sign bits of its DPLM-2
structure token. The transformer sees only this flat bitstream and its 1-D positions.

This representation separates what is shared from what is task-specific. Modality pairs may use
different frozen tokenizers, layouts, token budgets, model sizes, training regimes and inference
settings, but these choices introduce no modality-specific computation within a checkpoint's
denoiser.

\subsection{Partial denoising selects the task}
\label{sec:method:directions}

Let $M\in\{0,1\}^{S}$ mark the bits supplied as conditioning. We keep those positions clean,
corrupt the remaining positions, and evaluate \eqref{eq:loss} only on the generated subset:
\begin{equation}
\vx_\sigma=M\odot\vx_0+(1-M)\odot(\vx_0+\sigma\bm{\epsilon}),
\qquad
\Ls_M\propto\left\|(1-M)\odot(D_\theta-\vx_0)\right\|^2.
\label{eq:partial-denoising}
\end{equation}
The training regime samples masks for every supported direction. With no content clamped, the model
jointly generates all modalities. Clamping one segment gives cross-modal generation, and clamping a
prefix gives continuation. Thus, the task is specified by $M$, not by a task token, head or
encoder.

At inference, conditioned positions are clamped throughout the trajectory. Classifier-free
guidance combines a branch that sees the clean conditioning with one that sees the null value
$1/2$ at the same positions \citep{method:cfg}. Guided conditional sampling therefore uses two
network evaluations per integration step, whereas unguided joint generation uses one. Step count,
guidance and churn may vary by direction without retraining. One checkpoint consequently acts as a
joint generator and as every conditional generator included in its training regime.
